\chapter*{General introduction}
\addcontentsline{toc}{chapter}{General introduction}
\label{chap:introduccion-general-83}

% 00 — LA LINEA DEL HORIZONTE: ESCALA, FRONTERA Y MEDIDA
% ======================================================
% Frontera editorial sucesora; no importa contenido probatorio.
%
% Función: abrir la obra en la recta finita mediante intervalo, marca,
% hueco, cambio de escala, frontera de precisión y prolongación compatible;
% formular después el itinerario desde el estado discreto enriquecido hasta
% la publicación arquimediana.
%
% Declaración de continuidad: toda afirmación matemática conserva su
% fuente de REV.4 o su delta rector registrado. Este capítulo amplía la
% exposición principal; no sustituye ni abrevia las pruebas de referencia.

\ifdefined\HMTSucesoraStructureTrace
  \typeout{HMT-SUCESORA 00: opening of scale, boundary, and measure}
\fi

\section*{Inaugural thesis: the discrete structure of the continuum and the common origin of its publications}
\phantomsection
\label{sec:tesis-inaugural-helicoidal-hmt-md}
\addcontentsline{toc}{section}{Inaugural thesis: the discrete structure of the continuum and the common origin of its publications}

Triadic Modular Holography does not take the real line as a prior setting or
receive from it the values it will later use. Its input is finite: the
nine positions \(1,\ldots,9\), arranged horizontally and vertically so
that the same chart supports two irreducible operations. APP simultaneously preserves
the additive sheet, which accumulates, and the multiplicative sheet, which folds; the
residue, quotient and carry prevent equality between two published
values from erasing the difference between the histories that produced them. TRIT
orients each transition and distinguishes its regimes. TPK turns that
local orientation into dynamics: it selects the tritic mode without sacrificing
any APP sheet, transports, memorizes, returns and updates the state. Its
readings modulo \(3\), modulo \(9\) and modulo \(1000\), the nonadic connection,
the dodecaphase, the boundary and the return with memory are not independent
additions, but coordinated charts of the same genealogical operation.

Preservation is written from the first level:
\[
i+j=R^+_{ij}+9q^+_{ij},
\qquad
ij=R^\times_{ij}+9q^\times_{ij}.
\]
The residue publishes the class; the quotient unfolds its provenance. The
TPK transition preserves that distinction through
\[
\mathcal U_t=
\operatorname{Upd}_t\circ
\operatorname{Tra}_t\circ
\operatorname{Sel}_t,
\]
where selection incorporates the TRIT mode, transport realizes the
oriented direction and updating brings together memory and return. In the canonical
chart, the mode \(m_{\rm ph}=(1,-1,0)^3\) is common to the APP sheets and has
singular value \(s=1/2\); therefore
\[
1-s^2=\frac34=\frac{90}{120},
\qquad
\bigl\|[P_{\Sigma,3},P_{\Pi,3}]\bigr\|
=\frac{\sqrt3}{4}.
\]
The normalized commutator satisfies \(J_{\rm comm}^2=-I\). Sum--product
incompatibility, tritic orientation, defect \(90\to120\), complex
structure and electronic prefactor thus belong to the same internal
genealogy; they are not linked by decimal resemblance.

The compatible prolongation of these finite states produces increasingly
narrow cylinders and simultaneously preserves the history of their refinement.
The real line then appears as an Archimedean publication of the limit, not
as initial data. That chart first publishes the values of the nine
symbols of the seed and then prolongs the same genealogy toward
\(\pi\), \(e\), \(\varphi\) and \(\alpha\), without using any of them as a
retrospective selector. This is the thesis governing the entire work: the
joint discrete structure constructs the continuum and explains why its
stable values are what they are. The chain

\[
\{1,\ldots,9\}_{\mathrm h}\times\{1,\ldots,9\}_{\mathrm v}
\longrightarrow \mathrm{APP}\longrightarrow\mathrm{TRIT}
\longrightarrow\mathrm{TPK}\longrightarrow\mathcal X_{\mathrm{enr}}
\longrightarrow\mathcal S_\infty
\longrightarrow\mathbb R
\]

does not summarize heterogeneous chapters: it expresses a single mechanism seen at
successive resolutions. In it the phase may return while the state advances,
because the quotient, carry, orientation and memory have changed. The
cycle observed only through its phase is circular; with its genealogy preserved, it is
helical.

This difference is fixed through an injective lift. If
\(n=9q_9(n)+\rho_9(n)\), with \(1\leq\rho_9(n)\leq9\), the circular projection
\(\mathcal R_9\) and the helical immersion \(\mathcal H_9\) satisfy
\[
\mathcal R_9(n+9)=\mathcal R_9(n),\qquad
\mathcal H_9(n+9)=\mathcal H_9(n)+(0,0,1).
\]
The phase has closed and the genealogical height has grown: phase return and
state return are mathematically separated. The two-way involution
organizes conjugate ascending and descending sheets of the same lift, not
two independent generators. In this geometry, the helix preserves the
longitudinal advance of memory, the spiral records the transverse
contraction and the cylinder radius measures the width of the interval still
compatible with the prefix. The nesting of these half-open cylinders,
with radius tending to zero and genealogy preserved, is the mechanism that
subsequently publishes the real section.

The certified initial segment of the uniform prolongation constructing the
coinductive state \(\mathcal S_\infty\) is

\[
w_6\longrightarrow w_{12}\longrightarrow w_{18}\longrightarrow w_{24}
\longrightarrow w_{30}\longrightarrow R_{36}\longrightarrow G_9,
\]

but writing the chain acquires content only by declaring its
levels, thresholds and invariants. For
\(\chi\in\{\pi,e,\varphi\}\), let \(b_0^\chi=w_6^\chi\) and
\[
 b_{k+1}^\chi=b_k^\chi L_{\epsilon_k}\pmod 3,
 \qquad
 (\epsilon_0,\epsilon_1,\epsilon_2,\epsilon_3)=(0,0,1,1),
 \qquad
 w_{6(k+1)}^\chi=b_0^\chi\Vert\cdots\Vert b_k^\chi.
\]
The twelve transitions produced by \(\mathcal U_t\) form the matrix
pairs \(X_i,Y_i\) of the two regimes and uniquely determine
\(L_i=X_i^{-1}Y_i\), \(i=0,1\). The calendar
\((L_0,L_0,L_1,L_1)\) is the only one of the sixteen binary calendars that
simultaneously satisfies the three biographies; the one hundred and forty-four
single-entry perturbations fail. Therefore \(L_0\), \(L_1\) and
\(0011\) are rigid outputs of APP--TRIT--TPK transport, never data
chosen from \(\pi\), \(e\) or \(\varphi\). The matrices thus
constructed act on the current block; concatenation literally preserves
all preceding blocks. The visible word accompanies an enriched state that also transports
APP sheet, TRIT orientation, residue, quotient, carry, boundary,
memory and incidence shadow. Therefore, coincidence of words and
coincidence of states are distinct mathematical conditions.

\subsection*{Levels and thresholds of TPK prolongation}
\phantomsection
\label{sec:umbrales-prolongacion-tpk-inaugural}

\paragraph{\(w_6\): visible seed.}
It is the six-position ternary word published in each channel after
the chain of types
\(104\,976\to468\,U_6\to243\,w_6\subset\mathbb F_3^6\), with
\(|\mathbb F_3^6|=729\).
It fixes the first observable prefix, while the decimal emission \(U_6\), its
region, multiplicity and source state remain in the enriched
fiber. The action \(D_3\), the fiber invariants and the gauge
\((\operatorname{diag}_c,h_+)=(6,ES)\) fix the channel representative before
any comparison with its real value; in channel \(\pi\), five
distinct regions \(U_6\) also preserve different biographies over the
same word \(w_6^\pi\).

\paragraph{\(w_{12}\): first lift.}
It adds \(b_1=b_0L_0\) to the seed. It changes the visible depth from one to two
blocks and preserves \(w_6\) as an exact prefix, together with the channel and the
transported genealogy. The lift produces an extension of the state; it does not
replace the seed with a word without a past.

\paragraph{\(w_{18}\): closure of regime \(L_0\).}
It adds \(b_2=b_1L_0\). The second application of the same operator completes the
first calendar regime and distinguishes a recurrence from an isolated
coincidence: the eighteen trits contain, in their order, the levels of six and
twelve trits that precede them.

\paragraph{\(w_{24}\): typed change of lift.}
It adds \(b_3=b_2L_1\). Here the operator applied to the current block changes,
from \(L_0\) to \(L_1\), while the eighteen-trit prefix and the
transported state remain. The regime change modifies the continuation and preserves
the history on which it acts.

\paragraph{\(w_{30}\): first-boundary prefix.}
It adds \(b_4=b_3L_1\) and completes the finite calendar
\((L_0,L_0,L_1,L_1)\). Its five blocks per channel constitute the prefix
from which deep prolongability is tested. It is an entry level to
the boundary, while the same recurrence admits further blocks.

\paragraph{\(R_{36}\): first joint coinductive boundary.}
If \(u_5^\chi\) is the sixth block of channel \(\chi\), then
\[
 R_{36}:=(u_5^\pi,u_5^e,u_5^\varphi)^T,
 \qquad
 w_{36}^\chi=w_{30}^\chi\Vert u_5^\chi.
\]
The boundary relation selects that triple while preserving the axis
\(\varphi\), the lateral aggregate \(\pi/e\), saturation, column
occupancy and dual neutrality. Its equations leave exactly the two
mirror solutions
\[
 B_-=(021222,222220,102011)^T,
 \qquad
 B_+=(222220,021222,102011)^T;
\]
the APP orientation of the half-open cylinder fixes
\(R_{36}=B_+\). The subscript records thirty-six trits
exposed per channel; \(R_{36}\) is the boundary of the three sixth blocks and
\(\mathbb R\) is the subsequent Archimedean publication of the cylinder
system.

\paragraph{\(G_9\): connection of prolongation and return with memory.}
Level \(G_9\) organizes the nine local charts that act after
\(R_{36}\). If \(\mathcal R_K\) is the compatible extension relation in
phase \(K\), their composition defines the nonadic holonomy
\[
 \Gamma_{9,K}
 =\mathcal R_{K+8}\circ\cdots\circ\mathcal R_K:
 \widehat X_K\longrightarrow\widehat X_{K+9}.
\]
The subscript \(9\) comes from the phase cycle \(C_9\), inherited from the nine
APP positions, and counts the nine relations composing a turn. Its
operational content is certified by
\[
 g(K+9)=g(K),\qquad
 q_{{\rm mem},K+9}=q_{{\rm mem},K}+1,
 \qquad
 \widehat x_{K+9}\ne\widehat x_K.
\]
The phase label returns, memory advances and the prolonged state changes.
This is the structural meaning of \(G_9\): a connection turn over
an enriched fiber, whose cardinality comes from the architecture and whose
action is verified in transport. The modal quotient \(\mathcal K_{\rm ph}\)
occupies the subsequent reduced stage,
\[
 \mathcal U_t\longrightarrow\Gamma_9\longrightarrow\mathcal K_{\rm ph},
\]
and publishes compressed memory after the holonomy.
The local resolution of the ninth phase also distinguishes two horizons:
\(\#\operatorname{Surv}_1(B_9)=3\) and
\(\#\operatorname{Surv}_2(B_9)=1\). This descent \(3\to1\) certifies the
pruning of a specific fiber; composition of the nine relations certifies
the complete turn.

The dodecaphase records another coordinate of the same transport. The vertical
counter \(q_{{\rm mem},K}\) remains unbounded and its reading
\(\beta_K=[q_{{\rm mem},K}]_{12}\) satisfies
\(\beta_{K+9}=\beta_K+1\pmod{12}\). Reduction modulo twelve preserves the
dodecaphasic position; the integer counter preserves the number of turns. In
the terminal state, this memory feeds the signed record
\(U_{12}^{\rm sgn}\), the reversible seal \(K\) and the carry cocycle that
publishes the coordinate \(\alpha_{\rm HMT}\).
The nonsplit extension
\[
 0\longrightarrow C_{12}\longrightarrow C_{108}
 \longrightarrow C_9\longrightarrow0
\]
types the seam between the two calendars: nine steps return the phase and
advance one dodecaphasic sector; one hundred and eight steps close the observable clock
while the enriched state preserves the genealogical height reached.

Bidirectionality combines the channel mirror and the reversibility of the
lift within a single generator. The involution
\[
 \mathfrak c(u^\pi,u^e,u^\varphi)
 =(u^e,u^\pi,u^\varphi)
\]
exchanges the conjugate paths \(\pi/e\) and preserves the axis
\(u^\varphi\) and the aggregate \(u^\pi+u^e\). In the integer chart
\[
 x_kA=u_k+3x_{k+1},\qquad \det A=1,
\]
the advance determines residue and quotient, and the pair \((u_k,x_{k+1})\) recovers
\(x_k\). Ascent and descent are therefore conjugate paths of the same state;
descent reconstructs provenance without erasing accumulated memory.

Finally, \(G_9\) constitutes the recurrent unit of prolongation and not
its endpoint. The surviving histories form the inverse system
\[
 \cdots\longrightarrow\operatorname{Surv}_{\chi,n+1}
 \longrightarrow\operatorname{Surv}_{\chi,n}\longrightarrow\cdots,
 \qquad
 x_\chi\in\varprojlim_n\operatorname{Surv}_{\chi,n}.
\]
Nonemptiness, local uniqueness and compatibility of truncations
produce the coinductive section; contraction of its cylinders publishes a
unique real value, while the inverse limit preserves the genealogy that this
scalar value omits. The twenty blocks and the per-channel certificate of
\(396\) events, \(2376\) trits, \(43\) turns and \(387\) pairs
\((K,K+9)\) are finite executions
of that uniform law, not bounds on prolongation.

On the coinductive state thus constructed, a correlated kernel of
four publications is obtained. The correlation
expresses common origin and compatibility; it does not flatten their operational order. The
channels first publish the vectors \(P,E,\Phi\), while the signed terminal
reading constructs
\(U_{12}^{\rm sgn}\xleftrightarrow{\mathcal H_{12}}K\). On these outputs, and
without reintroducing their values as data, the dodecaphasic carry cocycle
derives the coinductive section \(\alpha_{\mathrm{HMT}}\), whose dodecaphasic
window is the rational projection \(\alpha_{12}\). The vacancy
kernel \(E_{21/22}\) provides a second internal construction: its
finite simple positive root \(\alpha_E\), prolonged by the same graded TPK
transport to the inverse system whose limiting root is
\(\alpha_B\). Integer closure defines
\(\alpha_A:=\varprojlim_N\rho_N(\alpha_{12})\), and the exact two-way
theorem establishes
\[
 \boxed{\alpha_A=\alpha_B=\alpha_{\mathrm{HMT}}}.
\]
The equality
\(\operatorname{Tr}_9(\alpha_E^{-1})=
\operatorname{Tr}_9(\alpha_{12}^{-1})\) is the first finite certificate of
that identity: it neither defines it, limits its scope nor selects any of
the prolongations. Thus,
\(\pi_{\mathrm{HMT}},e_{\mathrm{HMT}},\varphi_{\mathrm{HMT}}\) and
\(\alpha_{\mathrm{HMT}}\) belong to a single correlated coinductive generation,
but preserve differentiated stages, domains, readers and certificates. The
Machin identities and Leibniz brackets, the
Perron--Fibonacci realization, the factorial semigroup and the metrological tables subsequently recognize
or certify the publications already generated; none of them
selects the cylinders, vectors or carry of the generator.

The same common source has other publications that must not be confused
with one another either. In the internal mirror chart, the involution exchanges the branches
\(u^\pi\) and \(u^e\), preserves \(u^\varphi\) and the aggregate
\(u^\pi+u^e\); only after real publication are expressions such as
\(e+\pi\), \(e\pi\), \(\pi^e\) or the literal quotient \(\pi/e\) formed, each with
an arithmetic status that does not follow from the mere existence of the ternary
aggregate. In the incidence
chart, \(H^-_\alpha\), \(v_\alpha\) and \(N_\alpha\) denote,
respectively, a support, the vector of the exact index-three neighbor and the
resulting lattice; they do not denote the scalar \(\alpha_{\rm HMT}\) or intervene
in its numerical carry. The tetrad \(B_K\) enters the Witt star and the
vector \(v_\alpha\) determines the neighbor leading to
\(N_\alpha\cong\Lambda_{24}\). From the common Paley--Witt code \(C_W\),
two typed branches then open, not a single linear chain:
\[
 C_W\longrightarrow W_{12}\longrightarrow W_{24},
 \qquad
 C_W\longrightarrow N(A_2^{12})\xrightarrow{\,v_\alpha\,}\Lambda_{24}
 \longrightarrow V_\Lambda\longrightarrow V^\natural
 \longrightarrow\mathbb M\longrightarrow\text{Moonshine}.
\]
The scalar value of \(\alpha\) does not generate these symmetries. Numerical
derivation and exceptional prolongation are typed readings of the same enriched
state and preserve a common genealogical root.

Dimensional Mechanics begins when that structure is transduced into
action, magnitude, field and matter. From it come, according to its own
readers, the Planck constant and the electromagnetic,
thermal, gravitational and quantum constants; the mathematical structure of the electron,
the General Mass Law and mixing transports; and, later, the
Barbero--Immirzi parameter as a structural object of the transition between matter,
area and information. Conventional values appear at the end as
metrological recognition, never as inputs to the calculation. For the same
reason, resolution of the continuum problem and the five inseparable
constructions emerging from it precede terminal proofs concerning
classical mathematical problems: these stress-test an already
constructed structure, rather than generating it retrospectively.

The physical order is equally binding. The coordinate
\(\alpha_{\rm HMT}\) opens, in typed branches, the retention
\(\eta_{\rm ret}\) and the structural valuation \(10^{-34}\); their confluence
yields \(\hbar_{\rm ret}\), and only then are the angle and
counter-angle constructed and a single conversion between degrees and radians performed. The
electron brings this chain together in a material persistence. Its basal
\[
e_0=
\frac{\sqrt3}{4}
\left[
\exp\!\left(\frac{\varphi}{\pi^2}\right)
-22\alpha_{\rm HMT}^{\,3}
\right](1+15\Delta_4)
\]
\Needspace{8\baselineskip}
precedes beta closure; its two readers coincide at
\[
K_D(\Gamma_e)=K_\Omega(\Gamma_e)=(80,54,6),
\qquad
\kappa_e=80-54\alpha_{\rm HMT}+6\Delta_4,
\]
and the published mass is recognized metrologically only after evaluating
\(m_e^\beta=e_0(H_5/\eta_{\rm ret})^{\kappa_e}\). Tensor gravity, its
five-component algebraic sector and the typed physical realization; the
T-duality and M-theory screens; and the geometry of area and information
prolong the same antecedent with their own domains and falsifiers. None
of these realizations reintroduces as a premise the magnitude it must
explain.

Everything that follows develops this single thesis in different domains. Numbers,
symmetries and matter are not three juxtaposed collections. They are
three ways in which the constructed continuum publishes and returns the memory of
its own formation.

\section{Initial object and causal order: finite scale, observation horizon, and Archimedean publication}

\subsection*{Finite oriented chain and distinction between marks, cells, and genealogical units}

The exposition begins with a elementary operation: fixing to segment,
marking to scale on it and comparing the regions that esa scale determines. In
this point not it is presupposed yet the real line as complete support. The
initial object is a finite oriented chain
\[
 \mathcal I_K=(L_K,\prec_K;\,
 x_{K,0}\prec_K x_{K,1}\prec_K\cdots\prec_K x_{K,n_K};\,o_K),
\]
where \((L_K,\prec_K)\) is a finite totally ordered set, the
\(x_{K,j}\in L_K\) are observable marks, and \(o_K\) distinguishes the two
directions of traversal. No real magnitude has yet been assigned to \(L_K\).
The marks and consecutive cells are distinct objects: under the
closed-cell convention, the common mark is incident to the closures
of two adjacent cells without thereby becoming an additional cell.

This apparently minimal distinction determines the subsequent architecture.
A visible value describes the position published in a chart; a genealogical
unit further preserves the interval of provenance, the direction, the
accumulated quotient, and compatibility with finer scales. Therefore, two
numerically equal presentations may represent different histories,
and the same history may admit several coordinated publications.

For \(n\in\mathbb N\), the nonadic chart organizes this information by means of
the Euclidean division
\begin{equation}\label{p12-83-c0-s1:eq:horizonte-division-nonadica}
 n=9q_9(n)+\rho_9(n),
 \qquad q_9(n)=\left\lfloor\frac n9\right\rfloor,
 \qquad \rho_9(n)\in\{0,\ldots,8\}.
\end{equation}
The public phase numbering is obtained through
\[
 g(n)=1+\rho_9(n-1)\in\{1,\ldots,9\}\qquad(n\geq1).
\]
The residue localizes that phase, and the quotient \(q_9\) preserves how many revolutions
have been absorbed by the scale. This pair realizes a phase return
while preserving the trajectory. Its decisive invariant is the joint persistence
of both coordinates under refinement.

\subsection*{Inverse system of finite resolutions and relative observation horizon}

We shall call \emph{horizon} the boundary up to which a finite chart has
resolved the object and at which that description is linked with the
next chart. If
\(\mathcal S_K\) denotes the nonempty set of enriched states
compatible with the scale \(K\in\mathbb N\), the surjective links of
forgetting
\[
 \tau_{K+1,K}:\mathcal S_{K+1}\longrightarrow\mathcal S_K
\]
remove visible detail, but preserve the necessary information to
recognize a prolongation. We write
\[
 \tau_{K,K}=I,\qquad
 \tau_{L,K}\circ\tau_{M,L}=\tau_{M,K}
 \quad(K\leq L\leq M).
\]
The complete object is then the inverse limit,
\begin{equation}\label{p12-83-c0-s1:eq:horizonte-limite-inverso}
 X=\varprojlim_K(\mathcal S_K,\tau_{K+1,K}),
 \qquad
 x=(x_K)_K,\quad \tau_{K+1,K}(x_{K+1})=x_K.
\end{equation}
The nonemptiness of the levels and surjectivity of the links guarantee
the existence of compatible families; in the specialized channels, their
sections are subsequently constructed by means of the survival theorems.
Each truncation is finite; compatibility has no preset maximum
depth. Ten, one thousand, or ten thousand digits are checks of the same
prolongation rule, not alternative definitions of the number.

This formulation distinguishes three operations that decimal notation
usually superimposes. The first refines a partition. The second transports the
state between phases. The third publishes a coordinate in a continuous chart.
Only the last takes values in \(\mathbb R\). Thus, we begin with a finitely
marked line and end with the real line as the final Archimedean realization
of a compatible section.

\subsection*{Compatible augmentation and evolutionary conservation of the unit}

Unity is formulated before every Archimedean evaluation. Let
\(C_K=\mathbb Z[\mathcal S_K]\) be the free module of the states of level
\(K\). The link map induces
\((\tau_{K+1,K})_\#:C_{K+1}\to C_K\), and there exists the augmentation
\[
 \varepsilon_K:C_K\longrightarrow\mathbb Z,\qquad
 \varepsilon_K\!\left(\sum_{x\in\mathcal S_K}a_x[x]\right)=\sum_xa_x,
\]
compatible with the links,
\(\varepsilon_K\circ(\tau_{K+1,K})_\#=\varepsilon_{K+1}\). Evolutionary preservation
of unity expresses that refinement redistributes its content among
visible part, boundary, and memory without creating a second root or eliminating the
terminal term.

Let \(P_m\) be the cellular path with \(m\) vertices and \(m-1\) edges, and
\[
 Q_{m,d}=P_m^{\square d}.
\]
Its number of \(k\)-cells is
\begin{equation}\label{p12-83-c0-s1:eq:horizonte-celdas-cubicas}
 \#C_k(Q_{m,d})=\binom dk(m-1)^k m^{d-k}.
\end{equation}
The identity
\begin{equation}\label{p12-83-c0-s1:eq:horizonte-completacion-binomial}
 10^d=(9+1)^d
     =9^d+\sum_{j=0}^{d-1}\binom dj9^j
\end{equation}
is the equality of the \(0\) cells of the completion. The subcomplex
\(Q_{9,d}\) and its relative corona together form \(Q_{10,d}\). In dimension
three,
\[
 1000=729+271,
\]
where \(271\) counts only the new vertices. The vector cellular
relative complete is
\begin{equation}\label{p12-83-c0-s1:eq:horizonte-corona-relativa}
 \#C_\bullet(Q_{10,3},Q_{9,3})=(271,756,702,217),
 \qquad 271-756+702-217=0.
\end{equation}
Its edges, faces, and cubes cannot be omitted when studying incidence,
homotopy, or transport.

The adjective \emph{holographic} designates this compatibility here: the
boundary datum, accompanied by orientation, carry, and memory, controls the
prolongation of the state without replacing the interior by a number. The unit is
the same across dimensions, although its decomposition and readers
change.

\subsection*{Nonadic holonomy and phase return with memory increment}

Non-Abelian transport is a single connection over the phase cycle. Its
full holonomy is written
\begin{equation}\label{p12-83-c0-s1:eq:horizonte-conexion-nonadica}
 \Gamma_9=\operatorname{Hol}_{C_9}(\nabla_{\rm TPK}),
 \qquad g(k+9)=g(k).
\end{equation}
Here \(g:\mathbb Z\to C_9\), and for every point
\(\widehat x_k\) of an admissible orbit of the enriched state,
\[
 \Gamma_9\widehat x_k=\widehat x_{k+9},\qquad
 q_{\rm mem}(\widehat x_{k+9})
 =q_{\rm mem}(\widehat x_k)+1.
\]
In particular, when \(q_{\rm mem}\) distinguishes states,
\(\widehat x_{k+9}\neq\widehat x_k\).
The nine ordered components describe a turn of the connection; are not
nine filters applied to a motionless state. In each channel there exists a map
local typed
\[
 \operatorname{loc}_{K,\chi}:X_\chi\longrightarrow\mathbb F_3^6,
 \qquad |\mathbb F_3^6|=729.
\]
The transfer traverses the ambient fiber and locates the coordinate of the
compatible extension. The turn returns the phase and advances the memory,
the prefix, the carry, and the boundary.

In an isometric realization, the separation between visible publication and expressed memory is typed through
\[
 J:\mathcal H_{\rm vis}\hookrightarrow\mathcal H_{\rm full},
 \quad J^*J=I,\qquad
 U_{\Gamma_9}^*U_{\Gamma_9}=I,
\]
\[
 T=J^*U_{\Gamma_9}J,\qquad
 \eta=(I-JJ^*)U_{\Gamma_9}J.
\]
It follows
\begin{equation}\label{p12-83-c0-s1:eq:horizonte-compresion-expresion}
 U_{\Gamma_9}J=JT+\eta,
 \qquad J^*\eta=0,
 \qquad I-T^*T=\eta^*\eta.
\end{equation}
The term \(\eta\) expresses exactly the orthogonal defect produced by
this compression. For composable routes
\(\gamma_1:x_0\to x_1\) and \(\gamma_2:x_1\to x_2\), the carry
\(\operatorname{Carr}_R(\gamma)\in\mathsf C\) satisfies, with
\(H(\gamma),Y(\gamma)\in\operatorname{End}(\mathsf C)\), the cocycle law
\begin{equation}\label{p12-83-c0-s1:eq:horizonte-carry-cociclico}
 \operatorname{Carr}_R(\gamma_2\gamma_1)
 =\operatorname{Carr}_R(\gamma_2)H(\gamma_1)
  +Y(\gamma_2)\operatorname{Carr}_R(\gamma_1),
\end{equation}
which is the algebraic form of the genealogical continuity.

\subsection*{Dimensional elevation and coordination of geometric regimes}

The marked line provides the order and horizon. APP constructs then
the two-dimensional chart in which sum and product act on a support
common without losing their quotients. TRIT locally orients each interaction by
the unique residual decomposition
\[
 a+b=r+3c,
 \qquad a,b\in\mathbb Z,\quad
 r\in\{-1,0,+1\},\quad c\in\mathbb Z.
\]
The Topological Phase Kernel organizes routes, sheets, boundary, signature,
orientation, cylinder, record, memory, and transport. Elevation to the cube
does not add a spatial metaphor: it introduces the incidences and coherences that
make it possible to compare interior, crown, and publication in successive dimensions.

This sequence also prepares the articulation among three geometries. The
circular structure preserves phase; the complex structure represents the
quarter-turn on a real space \(V\) by an operator
\(\mathcal J:V\to V\) with \(\mathcal J^2=-I_V\); the hyperbolic structure
is realized in the Poincaré disk, whose boundary lies at infinite
hyperbolic distance. Their identifications will be effected by explicit
maps, not by equality of drawings. The geometric horizon thus anticipates
the thesis of the book: a published metric, the memory of what has been traversed, and the
openness of continuations are coordinated readings of an oriented
structure.

The double publication originates from a single source:
\[
 \operatorname{ev}^{\rm arq}_\infty:X_\chi\longrightarrow\mathbb R,
 \qquad
 \operatorname{ev}^{\rm inc}_\infty:X_\chi\longrightarrow\mathcal E_{\rm inc},
\]
where \(\mathcal E_{\rm inc}\) designates the limit of the boundary
incidence objects constructed in Part II. The first map publishes the
Archimedean value; the second preserves the exceptional genealogy.

\subsection*{Causal sequence APP–TRIT–TPK, continuum, and typed publications}

The following chapters construct APP, TRIT and the Topological Phase Kernel
from their finite data; develop the discrete structure of the continuum and
its two publications; derive the genealogies of the constants; integrate
the five inseparable realizations—spectral, solenoidal, gauge,
cohomological and determinantal—of the jointly constructed continuum; and then construct,
through Dimensional Mechanics, the physical realizations of the sectors that
have them. Archimedean
publication appears at the end of each chain, as a reader of an already
constructed object. This is the precise sense in which the work begins with a line
and ends in \(\mathbb R\).

\paragraph{Location of complete proofs.}
This opening fixes the objects and causal order. The complete proofs are found in the chapters devoted to the discrete measure complex, the APP gnomonic construction, TRIT oriented transport, the operator architecture of the Topological Phase Kernel, cubic completion, and double publication. Subsequent cross-references link each construction to its complete proof; the present opening fixes the common domain and the order of composition.


\section{Projective Formalization of Unity in Successive Scales and Dimensions}
\markboth{General introduction}{Projective Formalization of the Unit}
\label{p12-83-c0-s2:chap:ampliacion-conservacion-unidad}

\subsection{Oriented interval, marks, cells, and closing edge}

The construction begins with an oriented line and a declared scale. It does not
begin with the already completed set of real numbers, but with the
finite object
\begin{equation}\label{p12-83-c0-s2:eq:amp-recta-k}
 \mathcal I_m=\bigl([0,10^m],\,\mathcal M_m,\,\mathcal E_m,\,o_m\bigr),
 \qquad m\geq1,
\end{equation}
where \(\mathcal M_m\) is the set of marks, \(\mathcal E_m\) the set of
elementary intervals, and \(o_m\) an orientation. The formula brings together four
distinct objects. The segment provides the ordered support; the marks
publish positions; the intervals compare separations; the orientation
determines the direction of the reading. Suppressing any of these coordinates
changes the type of the object.

In the first chart, \([0,10]\), the visible nonadic positions are
\(1,\ldots,9\). Between them existen eight interior intervals. The step that
connects the last visible state with the first of the next turn it
is recorded as closure-and-carry edge. By eso the equality
\begin{equation}\label{p12-83-c0-s2:eq:amp-nueve-marcas-ocho-intervalos}
 9\ \text{visible marks}
 \quad\longleftrightarrow\quad
 8\ \text{interior intervals}+1\ \text{closure link}
\end{equation}
does not identify marks with intervals. It expresses how counting and measure
share a boundary without becoming the same operation.

The Euclidean division
\begin{equation}\label{p12-83-c0-s2:eq:amp-division-altura-fase}
 n=9q_9(n)+\rho_9(n),
 \qquad 0\leq\rho_9(n)\leq8,
\end{equation}
is now read in two directions. The residue \(\rho_9\) determines the published phase;
the quotient \(q_9\) preserves the accumulated height. When the phase
returns to its initial value after nine steps, the quotient has advanced.
The trajectory is helicoidal: one coordinate of the wheel has been recovered,
but not the complete state.

\paragraph{Algebraic interpretation of the residue–quotient decomposition.}
The equality \eqref{p12-83-c0-s2:eq:amp-division-altura-fase} is not merely an algorithm for computing remainders. It defines an extension: projection to the residue contracts the history, and the quotient preserves the information required to lift it. The elementary falsifier consists in imposing \(\rho_9(n+9)=\rho_9(n)\) and \(q_9(n+9)=q_9(n)\) simultaneously. The first equality is correct; the second contradicts Euclidean division.

\subsection{Projective system of scales and cubical completion to arbitrary depth}

Scaling is performed through refinement maps, not through a
graphic repetition of the same rule. For \(m\leq n\), let
\begin{equation}\label{p12-83-c0-s2:eq:amp-refinamiento-escalas}
 r_{n,m}:\mathcal I_n\longrightarrow\mathcal I_m,
 \qquad
 r_{m,m}=I,
 \qquad
 r_{n,m}r_{\ell,n}=r_{\ell,m}.
\end{equation}
The map \(r_{n,m}\) preserves the mark published at resolution \(10^{-m}\) and
contracts the detail lying below that threshold. The composition
equation ensures that forgetting two levels at once produces the same
result as forgetting them successively. This is the condition that makes it possible
to speak of a single history observed at different scales.

In three dimensions, the passage from the nonadic fiber to the complete decimal chart
is expressed by
\begin{equation}\label{p12-83-c0-s2:eq:amp-729-271-1000}
 1000=(9+1)^3=9^3+3\,9^2+3\,9+1
 =729+243+27+1=729+271.
\end{equation}
The equality contains a stratification. The \(729\) interior states
have no boundary coordinate; the next \(243\) states have one;
the \(27\) states have two; the apex has three. The number \(271\) counts the vertices
added by the decimal corona. It does not count the internal boundary of the
nonadic cube, whose cardinality is \(9^3-7^3=386\), nor the complete relative complex,
which includes edges, faces, and cubes.

For every dimension \(d\), let
\begin{equation}\label{p12-83-c0-s2:eq:amp-completacion-d}
 \overline C_9^{(d)}=(E_9\sqcup\{\star\})^d,
 \qquad \star\notin E_9,
\end{equation}
and denote by \(B_r(d)\) the stratum with exactly \(r\) coordinates equal to
\(\star\). Then
\begin{equation}\label{p12-83-c0-s2:eq:amp-estratos-binomial}
 |B_r(d)|=\binom dr9^{d-r},
 \qquad
 \sum_{r=0}^{d}|B_r(d)|=10^d.
\end{equation}
The equation describes growth of the scale without imposing a terminal
dimension. Each \(d\) has more states, but all belong to the same
completion law.

\subsection{Compatible product measure and conservation of the unit under projection}

The expansion of cardinals is accompanied by a conservation of measure. Assign
to each symbol of \(E_9\sqcup\{\star\}\) weight \(1/10\), and let \(\mu_d\) be the
product measure over \(\overline C_9^{(d)}\). If
\(\pi_d:\overline C_9^{(d)}\to\overline C_9^{(d-1)}\) forgets the last
coordinate, then
\begin{equation}\label{p12-83-c0-s2:eq:amp-conservacion-proyectiva}
 (\pi_d)_*\mu_d=\mu_{d-1},
 \qquad
 \mu_d\bigl(\overline C_9^{(d)}\bigr)=1.
\end{equation}
The first equality is the evolutionary conservation of unity. The dimension
increases, the number of states changes, and the normalized mass remains equal.
The boundary does not add a second unity: it redistributes the same unity among
compatible strata.

This conservation has to cellular version. For the path \(P_m\) and the
cubical product \(Q_{m,d}=P_m^{\square d}\),
\begin{equation}\label{p12-83-c0-s2:eq:amp-celdas-cubicales}
 \#C_k(Q_{m,d})=\binom dk(m-1)^k m^{d-k}.
\end{equation}
In particular, the relative crown of \(Q_{9,3}\subset Q_{10,3}\) has
cellular vector
\begin{equation}\label{p12-83-c0-s2:eq:amp-corona-completa}
 \#C_\bullet(Q_{10,3},Q_{9,3})=(271,756,702,217),
 \qquad 271-756+702-217=0.
\end{equation}
The null Euler characteristic certifies that the new vertices are
accompanied by the required incidences. Reducing the corona to \(271\) eliminates
precisely the transport that the theory must preserve.

\subsection{Geometric curvature and observation fibers: formal separation between metric and horizon}

Euclid's fifth postulate governs the global behavior of
parallels. In the Euclidean formulation, it fixes a geometry in which the sum of
the angles of a triangle and the distance between geodesics correspond to zero
curvature. Its modification opens elliptic and hyperbolic geometries.
The visible horizon of a geometric chart does not by itself constitute a
proof of the postulate: it is the boundary at which the chart ceases to distinguish
prolongations.

The correct relation between the two concepts is established through maps. A
geometry \((M,g)\) provides distances and geodesics; a remark
chart
\begin{equation}\label{p12-83-c0-s2:eq:amp-carta-horizonte}
 \operatorname{Obs}_{K}:(M,g)\longrightarrow Y_K
\end{equation}
publishes the accessible information to resolution \(K\). The relative horizon
of \(x\in M\) is the fiber
\begin{equation}\label{p12-83-c0-s2:eq:amp-fibra-horizonte}
 \mathsf{Hor}_K(x)=\operatorname{Obs}_K^{-1}
 \bigl(\operatorname{Obs}_K(x)\bigr).
\end{equation}
Curvature belongs to \(g\); indistinguishability belongs to the fiber
of \(\operatorname{Obs}_K\). The horizon can change when the chart is refined
without the curvature changing. This separation prevents identifying an
observer limitation with an intrinsic property of space.

In the Poincaré disk, the Euclidean boundary lies at infinite hyperbolic
distance. On the sphere, positive curvature closes the geodesics without an
analogous boundary. In the Euclidean chart, parallelism remains global. HMT
coordinates these regimes through the tritic parameter and the algebras
\begin{equation}\label{p12-83-c0-s2:eq:amp-tres-regimenes}
 \mathbb A_s=\mathbb R[J]/(J^2+sI),
 \qquad s\in\{+1,0,-1\},
\end{equation}
but preserves its metrics and codomains. The common algebraic unit does not erase
the difference geometric between \(\cos\), the parabolic branch and \(\cosh\).

\subsection{Refinement, survival and thresholds of compatible prolongation}

An indefinite expansion contains two orientations of the same scaling
operation. Toward precision, refinement distinguishes increasingly narrow
cylinders; toward extension, height records increasingly long revolutions.
Let \(\mathcal S_K\) be the set of states surviving at level
\(K\). The complete section is
\begin{equation}\label{p12-83-c0-s2:eq:amp-seccion-infinita}
 X_\chi=\varprojlim_K\mathcal S_K,
 \qquad
 x=(x_K)_K,
 \qquad \tau_{K+1,K}(x_{K+1})=x_K.
\end{equation}
Each \(x_K\) is finite; compatibility has no maximum depth.
An infinitesimal is not identified with a prolongation branch: the former
is a scale relation in a realization; the latter is a datum of the
inverse system. Only a reader that declares the metric and the evaluation can
transform the latter into the former.

The mathematics of thresholds appears when a prolongation changes regime
upon crossing a boundary. If
\(\operatorname{Surv}_K^{(N)}\) denotes the states of level \(K\) that admit
extension up to \(N\), then
\begin{equation}\label{p12-83-c0-s2:eq:amp-horizontes-supervivencia}
 \operatorname{Surv}_K^{(N+1)}\subseteq
 \operatorname{Surv}_K^{(N)},
 \qquad
 \operatorname{Surv}_K=\bigcap_{N\geq K}
 \operatorname{Surv}_K^{(N)}.
\end{equation}
The vacancy records the threshold at which a visible continuation ceases to
be compatible. It is not an unstructured void: it preserves the cause of
discarding, the carry, and the boundary that determine the next level.

\subsection{Archimedean evaluation of compatible sections and reconstruction of \(\mathbb R\)}

The Archimedean evaluation appears after constructing the section. For
a decimal chart of base \(10\), let
\begin{equation}\label{p12-83-c0-s2:eq:amp-evaluacion-arquimediana}
 \operatorname{ev}_{\mathbb R}(x)
 =a_0+\sum_{j\geq1}a_j10^{-j},
 \qquad a_j=\operatorname{Pub}_j(x_j).
\end{equation}
Convergence is obtained because the compatible cylinders have diameters
that tend to zero in the declared chart. The digit \(a_j\) is a coordinate;
the sequence \((x_j)_j\) is the history; the real number is the limit publication.
Therefore,
\begin{equation}\label{p12-83-c0-s2:eq:amp-cadena-recta-R}
 [0,10]\longrightarrow[0,1000]
 \longrightarrow[0,10^m]
 \longrightarrow X_\chi
 \xrightarrow{\operatorname{ev}_{\mathbb R}}\mathbb R
\end{equation}
is a type chain, not a naive inclusion chain.

The second publication of the same \(x\) preserves the exceptional incidence.
The work will traverse both directions and will finally return to
\(\mathbb R\). The return will then be a conclusion: the real line will have
been explained as an evaluation of the compatible state, and not used
as a premise that renders genealogy unnecessary.


\section{Historical genealogy of the continuum problem and HMT–MD mathematical formulation}
\markboth{Historical Genealogy of the Continuum}{Historical Genealogy of the Continuum}

\subsection*{The historical problem of the continuum, value and measure}

The continuum constitutes one of the foundational problems of mathematics. Arithmetic represents discrete units; geometry compares extensions; analysis completes sequences and constructs limits; measure theory assigns magnitudes to sets; physics converts those magnitudes into observables. The efficacy of these disciplines rests on the possibility of publishing a stable and reproducible result. The problem governing this work arises one level earlier: to determine what structure makes that stability possible and what information must be preserved in order to recover a value, a measure, and their symmetries from a common origin.

On the discrete line, a mark and the gap separating it from the next mark are
distinct objects. The mark corresponds to the operation of counting; the interval,
to the operation of measuring. The passage from \(0\) to \(10\) contains the nine
visible marks \(1,\ldots,9\), the eight interior intervals between
consecutive marks, and one lifted closing edge that carries the count into the
next turn. The modular reading preserves the phase; the quotient preserves
the attained height. When both coordinates are retained together, the
circular repetition becomes a helix, and the sequence of integers
preserves the history of its returns.

Triadic Modular Holography (HMT) develops this intuition as a
mathematical construction. A finite system of positions, sheets,
orientations, paths, and transports produces compatible families of states.
Its evaluations publish digits and values; its lifts preserve
residue, quotient, carry, boundary, and memory. The continuum is realized as
the image of a compatible section of the system of prolongations. Its
discrete structure resides in the genealogy that links all
resolutions.

This formulation distinguishes three levels. A \emph{digit} is a visible coordinate
produced by a reader. A \emph{value} is the image of a section
under a specified chart. A \emph{genealogical number} is the complete
section that preserves depth, operation, orientation, sheet, residue,
quotient, carry, route, boundary, memory, and survival. The decimal expansion
constitutes an evaluation of that object; compatibility among its
stages constitutes its continuity.

The same structure transforms the concept of symmetry. A symmetry is the
evolutionary preservation of identity under a transformation with memory.
A cycle may recover the visible position while simultaneously transporting the
state to another height, another sheet, or a new block. Invariance belongs to
the complete genealogy. The real value and exceptional symmetry therefore appear
as two coordinated projections of the same antecedent:
one contracts compatible cylinders, and the other preserves the incidence that the
Archimedean contraction leaves outside the published coordinate.

Dimensional Mechanics (MD) prolongs this construction into the physical domain.
The numerical section receives a spectral representation; persistence is
realized as a particle-route; the route ledger produces a signature; the
signature feeds a character and a hierarchy of towers; the dimensional chart
finally publishes a magnitude. The governing chain of the work is
\[
 \text{unit}\longrightarrow\text{state with memory}
 \longrightarrow\text{continuum, number and symmetry}
 \longrightarrow\text{particle, constant and quantity}.
\]

The scientific objective is unitary and is developed in five Parts. The formal core constructs APP, TRIT, and the Topological Phase Kernel; the discrete structure of the continuum unfolds its sections, the double projection, and the exceptional genealogy; the genealogy of constants produces \(\pi\), \(e\), \(\varphi\), \(\alpha\), the internal action scale, and its transducers; five specialization chains construct their internal antecedents and isolate the comparators leading to each classical recognition; Dimensional Mechanics then constructs the constitutive constants, mass, quantum statistics, information, and gravity. Each measurement is presented together with the object and operation that generate it; each metrological comparison is opened only after the internal output has been fixed.

The order of the work follows the constructive dependency: an object appears where
it begins to govern subsequent results. Each statement declares
its premises, its proof, and, where appropriate, the morphism of
physical realization. In this way, a subsequent metrological comparison does not
replace the genealogy that produced the output.

\section*{From the genealogical closure to the physical realization}

The thesis of this work attains four coordinated and irreducible closures:
genealogical, operatorial, matricial, and physical. It does not consist in adjoining a
set theory, a theory of operator algebras, or a quantum formalism
to a prior numerical construction. The four developments arise from
the same causal precedence: APP constructs the dual positional and
arithmetic realization; oriented TRIT distinguishes return, the torsional present, and opening;
TPK preserves value, sheet, orientation, residue, quotient, carry, boundary,
route, register, memory, survival, phase, and gauge. Only thereafter are
the compatible section, its readers, and its representations formed. The governing diagram
is, therefore,
\[
\begin{gathered}
 \mathrm{APP}
 \longrightarrow \mathrm{TRIT}_{\rm orientado}
 \longrightarrow \mathrm{TPK}
 \longrightarrow \mathcal X_\infty^{\rm cal},\\[2mm]
 \mathcal X_\infty^{\rm cal}
 \xrightarrow{\;\mathfrak R\;}\mathbb R,
 \qquad
 \mathcal X_\infty^{\rm cal}\times\mathcal C_{\rm exc}
 \xrightarrow{\;\mathfrak E\;}\mathcal E_{\rm exc},\\[2mm]
 C(X_m)\xrightarrow{\;\operatorname{diag}\;}M_{9^m}(\mathbb C)
 \xrightarrow{\;a\mapsto a\otimes I_9\;}
 \mathrm{UHF}(9^\infty)
 \xrightarrow{\;\pi_\tau(\,\cdot\,)''\;}R_{\rm hyp},\\[2mm]
 \mathsf P_{\rm CT}^{+}
 \xrightarrow{\;\mathfrak R_{\rm MD}^{\rm disc,+}\;}
 \mathbf{SpecRoute}\times\mathsf{Hilb}_{\mathbb R,J}
 \times\mathbf B\mathfrak M_D^+\times\mathbf B G_C.
\end{gathered}
\]
The four lines do not form a single linear chain. Real evaluation and exceptional projection are sister readers of the enriched section; the tracial branch passes through the cylindrical algebras and their unital inclusions; the physical realization begins with prepared routes endowed with the structures required by each component. This separation of domains prevents the von Neumann factor from being deduced directly from a section, the exceptional projection from being identified with an action on digits, or a physical magnitude from being presented as though it were a TPK state without a realization morphism.

\subsection*{Genealogical Closure: Extensionality, Infinity, and Decision}

The first closure determines what it means to preserve an object across all depths. Given an inverse system of finite spaces \((X_m,p_m)\), a section \(x=(x_m)_m\) is not characterized solely by the equality of its published values. The fiber-descent theorem establishes the conditions under which a finite reader descends to a quotient and distinguishes surjectivity, constancy on fibers, and uniqueness of the descended map. König's lemma guarantees the existence of a branch when the levels are finite, nonempty, and finitely branching; uniqueness further requires the complete character that selects a single compatible prolongation.

On a fixed section, the decreasing cylinders provide an exact
realization of the passage \(0\!-1\!-\!\infty\): the support contracts, the
height increases, the mass remains unitary, and the cylindrical measures
converge weak-* to \(\delta_x\). This result coordinates the finite Kronecker delta,
the conditional expectations between levels, and the Dirac delta as a point
measure, without identifying the latter with a differential operator. The
prefix foundation precludes infinite descending chains in each finite
presentation; coinduction controls compatible prolongation. The von Neumann
ordinals are lifted with a route-memory fiber, and cardinality appears as a
decategorification that does not retain the entire genealogy.

The same distinction requires cylindrical power, closed power, and
full power to be separated, as well as set-theoretic choice, natural choice,
algebraic choice, and effective choice. Encoding in ZF/ZFC preserves the
constructions and permits comparison with Extensionality, Foundation,
Power Set, and Choice; it does not make the real line or the set-theoretic universe
the origin of APP--TRIT--TPK\@. In the effective domain, decidability of each
finite window, existence of the limit, uniqueness of an instance, and the
impossibility of a uniform total decider are four distinct assertions.
The Gödel--Turing reduction affects only the fourth: it does not invalidate particular
selectors already constructed. Analogously, the Gödel--Cohen comparison separates
the extension \(M\mapsto M[G]\) from the internal enrichment of a state and
avoids conflating the crossing of depth-dense sets with genericity over a
model. Borel cylinders, comparison hierarchies, and the projective
minimax passage are introduced with their hypotheses of compactness,
continuity, and consistency, not as automatic consequences of inverse
notation.

This closure also has an informational formulation. If an update
\(U\) of the complete state is bijective, then
\(H(X\mid U(X))=0\), and the ideal logical Landauer term vanishes. If a
reader \(q\) preserves only the visible shadow, the discarded information is
localized in \(H(X\mid q(X))\). In the operatorial formulation, a
trace-preserving automorphism preserves entropy, whereas a trace-preserving
conditional expectation \(E_N:M\to N\) satisfies
\[
 S_\tau(E_Nh)-S_\tau(h)=D_\tau(h\Vert E_Nh)\geq0.
\]
The cost corresponds to forgetting the fiber, not to the reversible transport of the
enriched state; this identity does not assert null total physical dissipation nor
reduced computational complexity.

\subsection*{Operatorial Closure: From the Nonadic Unit to Continuous Dimension}

The second closure operationally realizes the evolutionary conservation of unity. For the complete system, one considers
\[
 A_m=M_{9^m}(\mathbb C),
 \qquad \iota_m(a)=a\otimes I_9,
 \qquad
 \tau_m(a)=9^{-m}\operatorname{Tr}(a).
\]
Identity \(\tau_{m+1}\circ\iota_m=\tau_m\) simultaneously preserves
unity and trace. The inductive limit is the UHF algebra of supernatural type
\(9^\infty=3^\infty\), and the weak closure of its tracial GNS representation is
the hyperfinite factor \(R_{\rm hyp}\) of type \(II_1\). Values
\(\operatorname{rank}(P)/9^m\) become dense, and the trace of projectors in
\(R_{\rm hyp}\) realizes the entire interval \([0,1]\): a continuous dimension
emerges from a tower of discrete inclusions preserving unity.

This construction is not to be conflated with the marked branch
\(j_{m,r}(a)=a\otimes p_r\). The isometries of an individual fiber generate in
norm the compact operators on the Hilbertian limit and, under weak closure, a
factor of type \(I_\infty\). Selection of one sector preserves the mark, but not the
unit of the tower; the simultaneous lifting of the nine fibers is what
produces the UHF--\(II_1\) branch. Nor may the UHF realization be identified
without proof with the crossed-product realization of the profinite clock: a
canonical conjugacy requires the action, measure, amenability,
ergodicity, and isotropy to be declared. The orthomodular logic of projectors, Boolean
contexts, normal states, and entropic normalization
\[
 S_\tau(9^m\rho)=S_{\rm vN}(\rho)-m\log9
\]
are derived within this operator construction and keep separate the
matrix state, its tracial density, and its statistical publication.

Topology, dimension, and composition do not collapse in the factor. The closure
transversal is recorded by means of
\[
 \mathfrak C_{\rm HMT}
 =\bigl(C(\Omega_\infty),R_{\rm hyp},\mathcal G_{\rm TPK}\bigr).
\]
The commutative spectrum preserves points and cylinders; the tracial factor preserves dimension and statistics; the groupoid preserves concatenation, admissible inversion, and the provenance of routes. This threefold architecture is formulated after the double projection and the exceptional corridor. It neither founds nor replaces them: Archimedean evaluation and the Paley--Witt--Golay--Leech--\(V^\natural\)--Monster--Moonshine chain remain the two coordinated readings of the same enriched antecedent.

\subsection*{Matrix Closure: Depth, External Size and Dilation}

The third closure introduces simultaneous control of three indices. If \(m\) measures genealogical depth, \(n\) the exterior size of a grid, and \(s\) the matrix level of its entries, the family \(\mathfrak M_{m,s}^{(n)}\) is bigraded by \((m,s)\) when \(n\) remains fixed and trigraded by \((m,n,s)\) when the exterior size also varies. In particular,
\[
 \mathcal K_{m,s}
 :=\operatorname{UCP}\!\bigl(C(X_m),M_s(\mathbb C)\bigr)
\]
has two independent operations: the genealogical restriction
\(\rho_m^s(\Phi)=\Phi\circ\jmath_m\) and the matrix compression
\[
 \operatorname{mc}_{\boldsymbol V}(\Phi_1,\ldots,\Phi_k)(f)
 =\sum_rV_r^*\Phi_r(f)V_r,
 \qquad \sum_rV_r^*V_r=I_s.
\]
Both commute exactly:
\[
 \rho_m^s\!\left(\operatorname{mc}_{\boldsymbol V}(\Phi_1,\ldots,\Phi_k)\right)
 =\operatorname{mc}_{\boldsymbol V}\!\left(
   \rho_m^{s_1}(\Phi_1),\ldots,\rho_m^{s_k}(\Phi_k)\right).
\]
The identity is functorial and distinguishes refinement of history of reduction from the observer.

Every depth-compatible family determines a unique UCP map
over \(C(X_\infty)\) and admits a common Stinespring dilation. In the
diagonal corridor, a single projective spectral measure \(P\) satisfies
\[
 E_x^{(m)}=V^*P([x]_m)V
\]
at every depth and on every cylinder. The minimal dilation is unique up to
unitary equivalence and exactly preserves what is seen by the reader; a direct sum
provides a faithful ambient realization that preserves all distinctions of the
algebra. On the UHF limit, a compatible family of UCP maps likewise admits a common
prolongation and dilation. The conditional expectation
\[
 \mathbb E_m=\operatorname{id}\otimes\tau_9:
 M_{9^{m+1}}(\mathbb C)\longrightarrow M_{9^m}(\mathbb C)
\]
preserves the trace and discards the final factor in a controlled manner; its Stinespring environment records the omitted digit. In general, no uniform finite auxiliary memory exists for all depths. Moreover, faithfulness of a representation does not by itself imply APP--TPK dynamical faithfulness: the latter requires an intertwiner compatible with route, monodromy, and composition.

The matrix theory is instantiated in verifiable objects. The qutrit quantum
magic squares with parameters \((9,3)\) and \((12,3)\) have, respectively,
81 and 144 effects, exactly unital row and column sums, and
declared simultaneous dilations; the same POVM systems produce quantum magic
cubes of orders nine and twelve. The order-twelve case is organized through
the fibration \(12\to4\times3\), with its gauge and its explicit
equivariance condition. By another route, the sum and product conditional expectations of APP
generate the algebra
\[
 C^*(P_{\Sigma,2},P_{\Pi,2})
 \cong\mathbb C^4\oplus M_2(\mathbb C)\oplus M_2(\mathbb C),
\]
and at every level \(s\) the central block realizes the free matrix interval
\([5I_s,9I_s]\). These constructions distinguish quantum Latin squares,
magic squares, projective measures, positive measures, and semiclassicality;
no identification is obtained merely because the objects share order
nine.

The common spectral antecedent of two readers is formulated only after
both have been constructed. For finite maps
\(R_m:X_m\to\mathcal A_m\) and \(Q_m:X_m\to\mathcal E_m\), the incidences
\[
 P_{a,e}=1_{\{R_m=a\}}1_{\{Q_m=e\}}
\]
refine the two diagonal spectral measures and recover their marginals. The
result does not prove joint injectivity or require one reader to be a
compression of the other. In the noncommutative regime, two mutually
unbiased qutrit PVM systems do not possess a sharp joint PVM; simultaneous realization requires
compatible noise or a larger environment. This difference in type preserves the
HMT double projection: the matrix theorem classifies realizations of readers
and does not rewrite their foundational antecedent.

\subsection*{Physical Closure: Common Domain, Two Orientations, and Localization}

The fourth closure constructs the domain from which a TPK route can simultaneously receive spectral, Hilbert-space, dimensional, and Cartan realizations. Let \(\mathsf D_{\rm sp}\), \(\mathsf D_Q\), \(\mathsf D_D\), and \(\mathsf D_C\) be the reference domains of those four representations, with faithful forgetful functors into the category \(\mathsf P_{\rm TPK}^{\rm enr}\) of prepared states and registered routes. The fiber product
\[
 \mathsf P_{\rm CT}^{+}
 :=\mathsf D_{\rm sp}
 \times_{\mathsf P_{\rm TPK}^{\rm enr}}\mathsf D_Q
 \times_{\mathsf P_{\rm TPK}^{\rm enr}}\mathsf D_D
 \times_{\mathsf P_{\rm TPK}^{\rm enr}}\mathsf D_C
\]
imposes that the four components have exactly the same antecedent.
Their canonical projections are functorial adapters and produce
\[
 \mathfrak R_{\rm MD}^{\rm disc,+}
 =\bigl(\mathfrak S_+,\mathcal Q_+,\mathfrak D_+,
        \rho_{C,+}^{\rm disc}\bigr),
\]
with values, respectively, in spectral routes, real Hilbert spaces
with complex structure, the positive dimensional monoid, and the discrete
Poincare group. The construction asserts the functor on the common domain; it does not
assert that every TPK route automatically admits all four lifts.

On a path \(\gamma=e_1\cdots e_r\), the total variation
\[
 V_+(\gamma)=\sum_jw(e_j)
\]
and the net oriented transport
\[
 V_{\rm or}(\gamma)=\sum_j\varepsilon(e_j)w(\underline e_j)
\]
are sibling readers. The loop \(ee^\vee\) proves that the former does not,
in general, factor through the latter: it has zero oriented transport and
positive variation. To obtain an invertible realization, the domain is restricted
to those edges whose spectral and Hilbertian components are invertible, whose
dimensional datum is an additive cochain, and whose Cartan cocycles are normalized
and antisymmetric. Only then does the localization
\(\Pi_\vee(\mathsf P_{{\rm CT},{\rm rev}}^+)\) produce the oriented functor
\(\mathfrak R_{\rm MD}^{\rm disc,or}\) and transport inversion component by
component. The positive reading is not eliminated, and the oriented reading is not
obtained merely by changing a sign.

This physical construction precedes the Planck and electron constructions
because it supplies them with a domain, representations, and transport rules, but
does not absorb their reference proofs or the General Mass Law. The
number--particle--route--enriched chronological record--signature--
character--towers chain preserves its causal direction; the absolute action scale,
the central electron, the mass families, CKM, PMNS, and the physical sectors
retain their own objects, domains, and conclusions. A metrological coincidence is
compared only after the internal output has been frozen. Extensions that still
require an intertwiner, a smooth limit, a proof of dynamical stability,
or a physical morphism are stated with those data as exact hypotheses. The
absence of any of those maps prevents the theorem from being extended to the
corresponding codomain, without altering the genealogical, operatorial, and
matricial theorems already proved in their domains.

The reading of the work follows, consequently, the order of dependence and not that of external recognition:
\[
\begin{aligned}
 \mathrm{APP}&\longrightarrow\mathrm{TRIT}_{\rm orientado}
 \longrightarrow\mathrm{TPK}\\
 &\longrightarrow\text{enriched section}
 \longrightarrow
 \begin{cases}
  \text{Archimedean evaluation},\\
  \text{exceptional projection},\\
  \text{cylinder algebras and tracial closure},\\
  \text{prepared routes and physical realization}.
 \end{cases}
\end{aligned}
\]
Each output preserves its domain, hypotheses, and precise obstructions. What is common is not
a forced equality between codomains, but provenance from the same
enriched state. This is the precise meaning of holographic closure: no
terminal publication exhausts the genealogy that produces it, and no subsequent
closure can replace the double direction --Archimedean and exceptional--
that preserves the unit with memory.


\subsection*{Historical perspective and development of the theory}

The history of mathematics can be read as a successive enlargement of the operations admitted on the unit. The integer made counting possible; the rational incorporated division; the real closed processes of approximation; the complex turned rotation into an algebraic operation. Each enlargement fixed a domain, an equivalence, and a representation. Equality published the result with enormous efficacy and at the same time absorbed much of the specific route that had produced it.

Geometry followed an analogous development. Euclidean construction fixed
incidence and measure; non-Euclidean geometries showed the coexistence
of formally coherent curvature regimes
\cite{EuclidElements,PoincareScienceHypothesis}. The trigonometric circle,
the complex plane, and the hyperbolic disk coordinate rotation, periodicity, and
transport through different state spaces. HMT places beneath these
realizations a common ternary state and constructs from it the
elliptic, parabolic, and hyperbolic branches.

Mechanics incorporated a relational dependence into the unit. The law of the
lever linked weight, arm, and equilibrium; Mach's critique of absolute
motion brought the relation among body, frame, and totality to the fore
\cite{ArchimedesEquilibrium,MachMechanics}; relativity formalized kinematic
invariances, local equivalence, and the geometry of the gravitational
field~\cite{Einstein1918Principles}; quantum theory introduced an
action scale and propagation through path amplitudes
\cite{PlanckNobelLecture,Feynman1948}. The HMT--MD program assembles these
requirements within a prior problem: constructing the state that preserves position,
orientation, operation, memory, and prolongation before publishing a number or
a magnitude.

The three terms of the name specify the method. \emph{Holography} denotes preservation of the genealogy and of the prolongation rule between resolutions. \emph{Modular} denotes the coordination of the chart modulo nine, the orientation modulo three, the residue--quotient readers, and decimal publication in base one thousand. \emph{Triadic} denotes both the local orientation \(\{-1,0,+1\}\) and the organization of the levels linking the finite support, the numerical section, and its physical realization.

\subsection*{Evolutionary Conservation and the Meaning of Symmetry}

HMT formulates symmetry as preservation of identity under transformation
with memory. An operation can recover the same visible phase without
recovering the same enriched state. If \(\Theta\) denotes the transport,
the relation
\[
 \operatorname{pr}_{\rm vis}\circ\Theta^9
 =\operatorname{pr}_{\rm vis}
\]
coexists with a nontrivial increment of carry or memory. Horizontal
repetition is the projection of a helical evolution. Identity is
preserved because its genealogy permits it to be recognized through change, not
because it remains motionless.

The extreme and mean holographic ratio realizes this principle in the cubic completion.
\[
 999=729+243+27,\qquad
 1000=729+243+27+1.
\]
The apex \(+1\) of the cubic completion and the terminal term \(+I\) of the
exponential closure are closure operations defined on distinct
domains. Their structural correspondence lies in the fact that both complete an
evolution and preserve the unit in the respective codomain. This conservation
organizes the nonadic return, the decimal carry, the inward growth of
the cylinders, and the outward growth of the incidences.

The dimensional prolongation is expressed through
\[
 10^d=(9+1)^d=\sum_{k=0}^{d}\binom dk9^k.
\]
After normalization by \(10^d\), the coefficients form a projectively compatible
family of measures whose total mass remains equal to one. The
interior, the boundary strata, and the apex redistribute that unit as
the rank changes. In the cubic realization
\(\mathcal C_9=(\mathbb Z/9\mathbb Z)^3\), the six oriented projections
preserve the APP support, and the face--vertex operator produces a space of
rank four. The signature \((3,1)\) emerges from that operator and provides the
finite antecedent of the subsequent spacetime reading. Dimension,
boundary, and conservation of unity are here results of the same system
of projections.

\subsection*{Algebraic Unification of Elliptic, Parabolic, and Hyperbolic Geometries}

The conceptual point of departure was to articulate three realizations that
mathematics ordinarily presents in distinct spaces: the complex plane,
the trigonometric circle, and the Poincare hyperbolic disk. HMT
unifies them through the family
\[
 \mathbb A_s=\mathbb R[X]/(X^2+s),
 \qquad s\in\{+1,0,-1\}.
\]
The elliptic branch contains an internal root of \(-1\) and the circular closure; the
parabolic branch contains the nilpotent threshold; the hyperbolic branch contains
the two idempotents and the Poincare opening. These three algebras arise
from TRIT and subsequently govern the sheets, the channels, the
Lorentz--Minkowski metric, and the physical reading of time.

The difference between sum and product provides the elementary operation that
articulates these geometries. Sum accumulates displacements and alternatives;
product composes stages and folds inputs onto a common class when
noninvertible elements intervene. APP preserves both evaluations on
a common support. Its spectral separation and the block projectors \(P_\pm\)
fix the intrinsic modes; the conditional projectors
\(P_{\Sigma,d}\) and \(P_{\Pi,d}\), together with their commutator, construct a
finite precursor of Heisenberg incompatibility. The
subsequent physical reading transforms that separation into the conjugate angular
coordinates \(A\) and \(C^*\),
oriented channels, and holonomy.

\subsection*{Construction of the formal kernel}

APP is constructed from the positions \(1,\ldots,9\) in two coordinates. The
kernel \(8\times8\), the publication of
the null class by means of \(9\) and the gnomon of seventeen cells completan the
chart \(9\times9\). The identification of opposite sides produces the torus
\[
 X_9=(\mathbb Z/9\mathbb Z)^2
\]
and its displacements generate the Cayley graph. The evaluations
\(\Sigma\) and \(\Pi\) act on the same vertices and paths.

TRIT is the ternary oriented state. In an elevation
\[
 a+b=r+3c,
\]
the residue \(r\) publishes the output and the quotient \(c\) preserves how it was
obtained. The triple differentiates orientation, threshold and memory. TPK is the
operational system of genealogical dynamics that incorporates the tritic selector
of the operating mode, transports the state and simultaneously preserves the two
APP sheets with their position, orientation, phase, carry, boundary, path,
record and survival. Its nonadic,
ternary and positional readers, its twelve-phase system and its lifts
form a single tree of domains and codomains. Mathematical dependencies
are the typed arrows of that tree; each map is presented with its domain,
codomain and a reproducible calculation.

With greater precision, APP is the quintuple
\[
 \mathcal A_9=
 \bigl(X_9,\Gamma_{\rm APP},\Sigma,\Pi,
       \operatorname{Path}(\Gamma_{\rm APP})\bigr),
 \qquad
 \Gamma_{\rm APP}=\operatorname{Cay}
 \bigl(X_9,\{\pm e_1,\pm e_2\}\bigr),
\]
with evaluations
\[
 \Sigma(i,j)=\operatorname{dr}_9(i+j),\qquad
 \Pi(i,j)=\operatorname{dr}_9(ij),
 \qquad
 \operatorname{dr}_9(n)=1+((n-1)\bmod9).
\]
The lift simultaneously preserves
\(i+j=R^+_{ij}+9Q^+_{ij}\) and
\(ij=R^\times_{ij}+9Q^\times_{ij}\); the published value and quotient
information belong to distinct coordinates of the same state.

The lift of TRIT, \(\widehat\tau=(r,c)\), preserves residue and carry. The fiber
over the visible symbol \(0\) contains \((0,-1)\), \((0,0)\), and
\((0,+1)\), whose prolongation trees may differ. This multiplicity
is the minimal form in which the same visible coordinate supports distinct
histories.

TPK is organized into eight operator classes before its particular
realizations are presented: support and state; internal selection; transport;
readers and quotients; memory and boundary; periodicity and return;
prolongation; and outputs. This classification separates the carrier objects,
the maps updating the state, the readers publishing its coordinates,
and the functors contracting the section toward an Archimedean,
exceptional, or classificatory codomain. Each TPK operator thus receives a domain, a
codomain, a rule, and its invariants before receiving an abbreviated name.

The nonadic phases are distributed among three operator fibers: \(\{1,4,7\}\) activates additive evaluation, \(\{2,5,8\}\) activates multiplicative evaluation, and \(\{3,6,9\}\) records the threshold. The transition \(9\to1\) restores the published phase and lifts the dodecaphase memory; the visible cycle is therefore the projection of a helix of enriched states.

Once the three objects have been presented, the work develops its consequences.
The central block
\[
 B_c=9P_++5P_-
\]
fixes the local exchange coefficient \(2\) and the spectral gap
\(9-5=4\); modular aggregation produces \(51-45=6\), and the unfolding over
the nine positions produces \(54=9\cdot6\). The mixed-signature form and
\(SO^+(1,1)\) construct the finite Lorentzian geometry. The cubic completion
constructs the extreme and mean holographic ratio, the discrete area law,
and the propagation boundary.

\subsection*{Uniform algorithm at finite depth and correlated coinductive
publication of \texorpdfstring{\(\pi,e,\varphi,\alpha\)}{pi, e, phi, alpha}}

Numerical generation is governed by the funnel
\[
 104\,976\to468\,U_6\to243\,w_6
 \to w_{30}\to R_{36}\to G_9.
\]
The channels of \(\pi\), \(e\), \(\varphi\) and \(\alpha\) preserve different typed
readers: angular closure, self-proportional propagation,
autoscale and signed dodecaphasic coupling. The internal relation
\(\operatorname{Ext}\), coinductive survival and the exhaustive
partition of the fiber produce at each level a unique prolongable cylinder
compatible with the preceding truncation; the nested family determines
a unique enriched section of every depth. Only afterward do the
Leibniz brackets, the factorial semigroup and Perron--Fibonacci locate
the already generated blocks Archimedeanly. The runs of one thousand and ten thousand
digits are extensive finite certificates of that uniform construction, not its
logical foundation.

The initial prolongation constructs
\[
 w_6\longrightarrow w_{12}\longrightarrow w_{18}
 \longrightarrow w_{24}\longrightarrow w_{30}
 \longrightarrow R_{36},
\]
and the internal census selects before the evaluation
\[
 w_\pi=010211,\qquad w_e^{(\mathrm{Euler})}=201101,\qquad
 w_\varphi=121200.
\]
The superscript specifies that \(w_e^{(\mathrm{Euler})}\) belongs to the reader of
Euler's number \(e\); it is an object distinct from the electronic path
\(\Gamma_e\) of Dimensional Mechanics.
For each prescribed finite depth, \(\operatorname{Ext}\), TPK
transport and compatible survival first produce exactly one
cylinder \(C_{\chi,n}\) among the \(729\) refinements. The rational reader
is then applied to the already generated cylinder: it constructs its internal bracket and
verifies \(j_{\min}=j_{\max}\) as an Archimedean localization certificate,
without selecting the output or supplying target digits to the generator. The
\cref{thm:generacion-monodromica-conjunta-rev7} proves, for
\(\chi\in\{\pi,e,\varphi\}\), nonemptiness, uniqueness and
compatibility of these survivors at every finite depth; therefore,
the partitions construct
\[
 c_\chi=(C_{\chi,n})_{n\ge0}
 \in\varprojlim_n\{C_{\chi,n}\}.
\]
The cylinders are nested, their diameters converge to zero, and the
EF--\(G_9\) lift preserves the enriched fiber. This separates
the uniform algorithm parameterized by depth, the coinductive section it
determines, and its finite runs of one thousand and more than ten
thousand digits, which act as extensive certificates and not as the logical
foundation of the theorem.

The three sections realize distinct structural functions. \(\pi\)
organizes closure, orientation, and monodromic boundary; \(e\) organizes
self-proportional propagation and prefix evolution; \(\varphi\) organizes
Perron--Fibonacci self-scaling and its stable mode \(\varphi^{-2}\). The
oriented operator \(\mathscr R(x,y)=x/y^2\) publishes the conjugate readings
\(\pi/\varphi^2\) and \(\varphi/\pi^2\), which reappear in the physical
areal--volumetric hinge.

The terminal state supplies the block and signed record that feed
the dodecaphasic closure of \(\alpha_{\rm HMT}\). The carry recurrence in
base one thousand and the equation \(E_{21/22}\) constitute two internal pathways with
distinct sources. Next, number theory organizes the
complementary constants through extraction operators: finite parts,
Laurent jets, \(L\) functions, prime products, Gauss dynamics,
Rogers--Ramanujan and Feigenbaum renormalization. The text distinguishes
a common antecedent from equality of value; for example, \(\varphi\),
\(\varphi^{-1}\) and \(\varphi^2\) share an algebraic genealogy without
becoming the same number.

The complete genealogy of the first path is
\[
 \begin{aligned}
 104\,976&\to468\,U_6\to243\,w_6\to w_{30}\to R_{36}\to G_9,\\
 G_9&\to x_{\rm term}\to(B_{\rm term},U_{12}^{\rm sgn})\to K.
 \end{aligned}
\]
The seal \(K\) and the twelve terminal triads feed the carry
\[
 \begin{aligned}
 s_m&=P_m+E_m-\Phi_m-K_m+c_{m+1},\\
 A_m&=s_m\bmod1000,
 &c_m&=\frac{s_m-A_m}{1000}.
 \end{aligned}
\]
with \(c_{13}=0\). The resulting integral identity fixes
\[
 \alpha_{12}=\pi_{12}(\alpha_{\rm HMT})
 =\frac{7297352569283800997285105472380663}{10^{36}}.
\]
This identity publishes the exact thirty-six-digit rational window;
the compatible recurrence at arbitrary depth constructs
\(\alpha_{\rm HMT}=\varprojlim_N A^{(N)}\), which is not reduced to a single
window. The second pathway solves the vacancy equation \(E_{21/22}\)
through an analytic kernel with a simple positive root \(\alpha_E\). The
graded TPK transport prolongs that kernel into a compatible family
whose limiting root is \(\alpha_B\); integer closure determines
\(\alpha_A:=\varprojlim_N\rho_N(\alpha_{12})\). The truncation squares
of both families commute and determine the same surviving cylinder at
each depth, so
\[
 \boxed{\alpha_A=\alpha_B=\alpha_{\rm HMT}}.
\]
The equality after \(\operatorname{Tr}_9\) is only the first finite
certificate of this exact identity; it does not constitute its criterion, bound its
extension or select the prolongations. Each pathway preserves its own boundary
condition within the same enriched state.

Number theory extends the same criterion of object, operator, and
extraction. Euler--Mascheroni appears as a finite part; the Stieltjes constants
as Laurent coefficients; Apery as the evaluation of
\(\mathcal Z_3(3)\); Glaisher--Kinkelin through regularized derivation;
Euler--Kronecker through an Eisenstein character; Meissel--Mertens through a
prime trace; Catalan through \(L(2,\chi_{-4})\); and Khinchin through
Gauss dynamics. Rogers--Ramanujan, the prime clocks, and the Feigenbaum
approximants belong to distinct operator constructions. The
combinations \(e+\pi\), \(e\pi\), and \(\pi^e\) are formed after their
arguments have been constructed and preserve the arithmetic domain
corresponding to each statement.

The radical \((3)\subset\mathbb Z/9\mathbb Z\) organizes the visible network
\(3\)--\(6\)--\(9\) and the orbit \(369\to693\to936\). Over the same
arithmetic, the power--vacancy theorem
\[
 V_9(n)=\left\lceil n\log_{10}\!\left(\frac{10}{9}\right)\right\rceil
\]
gives the exact evaluation
\[
 \operatorname{ord}_{10}\!\left(9^{9^9}\right)=369\,693\,099,
 \qquad
 \#\operatorname{cifras}\!\left(9^{9^9}\right)=369\,693\,100.
\]
This figure is not the census of \(396\) elementary boundary events per
channel: the former measures the
decimal order of a nonadic power; the latter counts states in an
emission calendar.

\subsection*{Double projection as the sole thesis}

The real continuum is the Archimedean shadow of an enriched section. The
projection \(\mathcal P_{\rm arq}\) contracts compatible cylinders and publishes
the value. The projection \(\mathcal P_{\rm exc}\) preserves support,
orientation, carry, and memory. Both proceed from the same state; therefore,
number and exceptional symmetry form a single causal thesis.

The incidence branch constructs Paley, Witt, the ternary Golay code and
the Steiner systems. From the common code \(C_W\), the lift
\(W_{12}\to W_{24}\) remains distinct of the branch
\[
 C_W\to N(A_2^{12})\to\Lambda_{24}\to
 V_\Lambda\to V^\natural\to\mathbb M,
 \qquad [g]\longmapsto T_g.
\]
The configuration of the twenty-seven lines, the Schläfli graph, and
\(W(E_6)\) receive their own incidence operators. The definition of the
lattice VOA, the orbifold involution, and the replicable functions show
how genealogical identity is preserved when changing codomain. This
structure also provides the rank screens and the mathematical
interface with T-duality, strings, and M-theory.

Level-five modular arithmetic enters through the
Rogers--Ramanujan continued fraction \(R(q)\). The Klein--Ramanujan identity expresses
\(j(\tau)\) rationally in \(R(q)^5\), and the limit
\(R(q)\to\varphi^{-1}\) links pentadic self-scaling with the modular branch.
On the canonical branch of order two,
\[
 t_2=q^{-1}\prod_{n\ge1}(1+q^n)^{-24},\qquad
 T_{2A}=t_2+24+\frac{2^{12}}{t_2},
 \qquad j=\frac{(t_2+256)^3}{t_2^2}.
\]
The lattice VOA \(V_\Lambda=M(1)\otimes\mathbb C[\Lambda]\), the involution \(\theta=-I\), and the fixed--twisted orbifold realize the lattice output in the domain of vertex operators. The classical FLM identification fixes \(V^\natural\) and \(\operatorname{Aut}(V^\natural)=\mathbb M\). HMT supplies the antecedent and the incidence maps that reach that construction; the VOA identification theorems retain their own mathematical domain.

The finite root and the real structure are coordinated through the integral presentation.
\[
 \mathcal Q_{\mathbb Z}=\mathbb Z[u]/(u^2+1),
\]
not by means of an inclusion between fields of different characteristics. The
Witt endomorphism \(A_W^2=-I_6\) and the real operator
\(J_{\mathrm e}^2=-I\) are specialized representations of that relation.
The quarter-turn \(z\mapsto\mathrm iz\) preserves both the Euclidean metric
and the Poincare metric of the disk. The identity
\[
 \operatorname{Exp}(\pi_{\rm HMT}J_{\mathrm e})+I=0
\]
is applied only after the Archimedean reader has generated
\(\mathcal P_{\rm arq}(x_\pi)=\pi_{\rm HMT}\). The exceptional branch supplies
the quadratic type, not the angle: value and incidence remain coordinated without
either branch acting as an oracle for the other.
The family \(J_s^2=-sI\),
\(s\in\{+1,0,-1\}\), extends Euler to the elliptic,
parabolic, and hyperbolic realizations. The final term \(+1\) and the apex of EMRH
perform closure functions in distinct and coordinated domains. This
relation provides an operator expression of the evolutionary preservation of
the unit without identifying objects of different types.

\subsection*{Planck scale, Barbero--Immirzi functional, and CKM--CP mixing}

The double projection causally precedes the angular and physical
constructions. The Smith normal form of \(H_4\) produces a cokernel
of order sixteen and the transfer
\[
 \Sigma_H=\varphi\alpha_{\rm HMT}^{16}
\]
proves
\(\lfloor\log_{10}\Sigma_H\rfloor=-34\) and thereby fixes the internal scale
\(10^{-34}\mathcal U_S\). The action line contains two coordinates,
\[
 \hbar_{\rm pre}=H_5\,10^{-34}\mathcal U_S,\qquad
 \hbar_{\rm ret}=\eta_{\rm ret}\,10^{-34}\mathcal U_S.
\]
Temporal periodicity of order \(108\) supplies a self-dual
Planck closure. In the circular realization, with
\(\ell_0:=c_{\rm int}t_0\),
\[
 T_{108}=108t_0,
 \qquad
 R_{108}=\frac{54}{\pi}\ell_0,
 \qquad
 \theta_{108}^{\rm circ}=\left(\frac{54}{\pi}\right)^2,
\]
and the quantum satisfies
\[
 R_{108}=\bar\lambda_{\rm C}=r_{\rm g}=\ell_{\rm P}^{\rm circ}.
\]
The broken line
\[
 Z_a(m)=\frac{m_{{\rm P},a}^{\circ}}mP_+
       +\frac m{m_{{\rm P},a}^{\circ}}P_-,
 \qquad N_{\rm sp}(Z_a(m))=1,
\]
interchanges Compton and gravity under \(m\mapsto(m_{{\rm P},a}^{\circ})^2/m\). Its fixed point is the Planck scale. The electron section constitutes a subsequent genealogical realization on that line and not an identification of the electron with a classical black hole.
The direct angular coordinate \(A\) and its conjugate angular coordinate \(C^*\) subsequently construct the channels \(q_\pm\). The hexad--octad inclusion produces the Barbero--Immirzi functional, the area operator, and its spectrum. The same rank-three incidence constructs CKM, the CP phase, and the Jarlskog invariant. These results belong to the mathematical closure of HMT; their field interpretations are resumed in MD.

The Smith normal form explicitly fixes
\[
 \operatorname{SNF}(H_4)=\operatorname{diag}(1,2,2,4),
 \qquad |\operatorname{coker}H_4|=16,
\]
and the action line distinguishes
\(\hbar_{\rm pre}\), \(\hbar_{\rm ret}\) and
\(h_{\rm ret}=2\pi\hbar_{\rm ret}\). The chart
\(\mathcal U_S\mapsto\mathrm{J\,s}\) publishes then the unit
metrological; the order \(-34\) belongs to the internal section that precedes to
esa chart.

The combinatorial inclusion of a hexad into its associated octad produces
\[
 6\binom62=90,\qquad
 8\binom62=120,
 \qquad
 \frac{90-120}{24\binom62}=-\frac1{12}.
\]
On these spaces, the dodecaphasic fundamental chain brings together twelve sectors
and a single oriented return seam. Its linear reader produces
\(12S_{90}-S_{120}\) and, therefore, the weights \(12:-1\) of the functional
\(\Gamma_{6\to8}\). The coefficient of \((1-q)^{-1}\), equal to \(1/24\),
positivity and minimality constitute subsequent arithmetic
checks of the pair produced. The Holst specialization
\(\gamma=\Gamma_{6\to8}\) is realized at the triadic boundary through
\[
 \widehat A_{\rm HMT}
 =\gamma_{\rm HMT}a_0\sum_{e\in\partial A}N_e,
 \qquad
 \operatorname{Spec}(\widehat A_{\rm HMT})
 =\{n\gamma_{\rm HMT}a_0:0\le n\le72\}.
\]
The collective intertwiner satisfies
\(J_{6\to8}D_{\rm inc}=D_{\rm area}J_{6\to8}\) on the collective incidence
subspaces. The area operator, its spectrum, and the triadic,
bilayer, and modular entropies thus constitute the internal closure that realizes the
Barbero--Immirzi parameter. The rank-three rational transformation
constructs, in a sister branch, the CKM angles, the CP phase, and the Jarlskog
invariant before experimental comparison.

\subsection*{Genealogical Closure of the Number and Physical-Spectral Realization}

The first four Parts culminate where a conventional exposition usually begins:
they define number as a compatible section and value as its projection.
Dimensional Mechanics preserves that definition and adopts a
physical--spectral reader. It first
reconstructs the APP--TRIT--TPK kernel in physical--spectral form. APP is realized
as a two-dimensional torus and distinguishes additive, multiplicative, areal,
volumetric, and torsional sheets. TRIT receives the semantics
\[
 \begin{aligned}
 +1&=\text{return, memory and past},\\
 0&=\text{threshold and present},\\
 -1&=\text{opening and compatible future}.
 \end{aligned}
\]
TPK transports the route, the record, and the holonomy. Exact periodicity
constructs the finite dynamics; a physical time crystal adds a Hamiltonian,
robust subharmonicity, and stability. The conjugate angular coordinates
\(A\) and \(C^*\) spectrally realize the sum--product difference.

The gnomonic completion here acquires an explicit geometric function. The L carries the Euclidean chart onto which the inertial and gravitational readers project to a unit quotient. This chart constitutes the realization of the local equivalence principle. The enriched state prior to projection keeps the areal and volumetric sheets separate, and the Ambivalence Principle fixes their ratio \(\pi_{\rm HMT}/\varphi_{\rm HMT}^{2}\). Together with the negative- and positive-curvature charts, this yields an AdS--Euclidean--dS atlas: the first realizes the correspondence between interior and boundary, the second governs local physics, and the third satisfies \(\ddot a/a=H^2>0\) and geometrically explains accelerated expansion. The three charts coexist because they are realizations of the same TRIT triad, linked by TPK transport.

The particle is defined on this grammar. A particle-route is a persistent
section endowed with orientation, sheet, occupation, action, propagator,
record, and anti-section. The mass law is first formulated on each
multisection \(F\):
\[
 \widehat{\mathbf M}^{\beta,r}_F
 =\mathbf m_\star^{\rm ret}\otimes
  R_{\rm act}^{\widehat K_F^r/2}
  \widehat{\mathcal A}^{(0)}_F
  R_{\rm act}^{\widehat K_F^r/2},
 \qquad r\in\{D,\Omega\},
\]
where
\[
 \mathbf m_\star^{\rm ret}
 =\eta_{\rm ret}\,10^{-34}\mathcal U_S
   \frac{2\pi}{108t_0}\,c_{\rm int}^{-2}
\]
fixes the absolute mass section and
\(\widehat{\mathcal A}^{(0)}_F\) gathers signature, character, and all the towers of
\(\langle90,120\rangle\). The complete atlas is thus the domain of the law; each
species subsequently selects its realization within the corresponding fiber.

The central electron is the first rank-one scalar reduction of this
definition. The global torsion entering its reader is
\[
 \Delta_4=\frac{\pi-e\log\pi}{270}
 \left(1+\frac{\pi}{729}\right).
\]
Its bidirectional closure produces
\[
 \begin{aligned}
 R_{\rm act}&=\frac{H_5}{\eta_{\rm ret}},
 &\kappa_e&=80-54\alpha_{\rm HMT}+6\Delta_4,\\
 \mathfrak e_e^\beta&=\mathfrak e_e^{(0)}R_{\rm act}^{\kappa_e},\\
 \mathfrak e_e^\beta
 &=0.5109989506900008086082571648\ldots\ {\rm MeV},
 &m_e^\beta&=\mathfrak e_e^\beta/c^2.
 \end{aligned}
\]
The identity
\(\Sigma_H=\varphi\alpha_{\rm HMT}^{16}\) simultaneously supplies the
scale \(10^{-34}\mathcal U_S\) to the two action coordinates; the
quotient \(R_{\rm act}\) preserves the relative refinement and cancels that
common scale. Thus, the internal action scale, the bidirectional correction, and the
energy chart in MeV occupy successive levels of a single derivation.
The first quantity is the published mass energy; the second is the mass
obtained through the explicit chart $E=mc^2$. This distinction precludes identifying
an energy coordinate with a unit of mass.
Vacancies construct polarization and current; the response of the sheets
constructs the permittivity and permeability of the vacuum. The SI chart publishes
\[
 \begin{aligned}
 \mu_0^{\rm HMT}
 &=1.256637062121047789\ldots\times10^{-6}\,\mathrm{H\,m^{-1}},\\
 \varepsilon_0^{\rm HMT}
 &=8.854187812793002329\ldots\times10^{-12}\,\mathrm{F\,m^{-1}}.
 \end{aligned}
\]
The atlas then situates
the quarks and remaining species, defining color, flavor, charge, spin,
statistics, and representation within the same genealogy.

In this classification, \emph{colour} is the ternary fibre that organizes the
three realizations of a quark route and whose neutral composition satisfies
\(0+1+2=0\pmod3\). \emph{Flavour} is the coordinate that distinguishes, within
the same representation, the sections transported by the mixing
operator. \emph{Spin} labels the representation of the cover of the
rotation group carried by the section. The holonomy of the oriented band
realizes its closures \(2\pi/4\pi\) and distinguishes phase return from orientation
return. Statistics follows from the
grading and from the symmetric or exterior realizations; once CCR or CAR is fixed, it
publishes Bose--Einstein or Fermi--Dirac occupation and, in the latter case, the
Pauli exclusion principle. A virtual particle is a finite prolongation that
conserves route, record, and signature up to a declared horizon; a prolongable
material section conserves the corresponding global compatibility.

\subsection*{Operational domain of paths, fibers, and physical realizations}

The construction proves that the internal genealogy reaches the complete
physical domain. The domain distinguishes
\(81\) positions, \(56\) families, \(13\) multisections, and \(324\) oriented
routes, together with the null sector. Each realization preserves, according to its codomain,
position or fiber, family, route, orientation, phase word, register,
signature, character, refinement operators, chart, and comparison. External
values are joined by identifier after the internal output has been frozen.

The construction chain that governs those rows is
\[
 \begin{aligned}
 \text{surviving section}&\to\text{observable route}
 \to\text{record}\to\sigma\in\mathbb Z^5,\\
 \sigma&\to\chi_0\to\chi_{\rm full}\to\text{towers}
 \to\text{dimensional section}.
 \end{aligned}
\]
The two bidirectional readers publish
\[
 K_D=\left(D_{27}+D_{36},D_1,\frac{D_4-D_3}{2}\right),
 \qquad
 K_\Omega=(\Omega,54,P_6),
\]
and act as diagonal operators on each fiber. In the central cell, both reduce to \((80,54,6)\); in the noncentral fibers, their spectra preserve the distinct routes before the selection of a physical realization. The operatorial comparison presents first the HMT output, then the chart, and finally the Standard Model or PDG reference with its unit, scheme, uncertainty, and residual.

The scientific output assumes the type corresponding to each sector. The electron
cell publishes the closed scalar beta mass; the photon and gluons publish the
null chart; the noncentral fibres preserve their operators and spectra.
The finite flavour--generation--colour morphism to a persistent route is already
constructed, and the twelve nominal routes
\(e,\mu,\tau,u,d,s,c,b,t,W,Z,H\) publish their angular signatures and their
scalar characters before external comparison. In leptons and quarks, the
operators \(K_D/K_\Omega\) also refine the fibre; they are not the requirement that
retroactively creates its scalarization. \(W\), \(Z\), and \(H\) preserve their
native CT--108 reader even though they do not form part of the operatorial atlas of thirteen
multisections. The topological sectors publish fusion and braiding; the
virtual sections publish their prolongation horizon. This
heterogeneity preserves the proper codomain of each class and makes the
table a complete comparison of typed results, values, and
realization boundaries.

\subsection*{Quantum, information, gravity and cosmology}

The propagation of paths constructs the finite sum over paths, the
Clifford--Dirac representation and the Schrodinger evolution. Fixed the
Hamiltonian of the internal calendar, the evolution satisfies
\[
 i\hbar_{\rm int}\dot\psi=H_{\rm clk}\psi,
\]
and its finite propagator admits the expansion
\[
 (U^R)_{fi}=
 \sum_{\gamma:i\to f,\,|\gamma|=R}
 \prod_{\ell\in\gamma}U_\ell.
\]
The product preserves the order of transitions within a history, and
the sum accumulates alternative histories with common endpoints. The grading
of the exchange is realized concretely on
\(\mathbb T_{27}^{\,2}\times\{\mathrm N,\mathrm S,\mathrm E,\mathrm W\}\)
through
\[
 U_{\rm micro}=S D_J C_4,
 \qquad U_{12}^{\rm reg}=U_{\rm micro}^{12}.
\]
This propagator is unitary; its powers enumerate the finite routes exactly, and its discrete transform produces \(27^2\) Fourier blocks. The construction is not thereby identified with a continuous functional integral. The grading fixes the exchange datum; the representation of the permutation group determines the symmetric or exterior completion. Once the CCR or CAR relations are given, the variational principle and entropy maximization produce the Bose--Einstein and Fermi--Dirac distributions, and the latter realizes Pauli exclusion. In the Möbius basis, one visible turn of \(2\pi\) changes sheet and accumulates \(h/2\); the double visible-phase return of \(4\pi\) restores orientation and accumulates \(h\). Per radian of internal phase, the transport remains equal to \(\hbar\).

The Boltzmann constant appears as a transducer between action and ternary
information. The internal construction fixes the product
\(k_B\Theta_{\rm clk}\) and its structural function; the BIPM definition of the kelvin
then fixes the numerical coordinate of \(k_B\) in the SI chart:
\[
 k_B\Theta_{\rm clk}\log3=\hbar_{\rm int}\frac{2\pi}{108t_0},
 \qquad
 k_B^{\rm SI}=1.380649\times10^{-23}\ \mathrm{J\,K^{-1}}.
\]
The ensembles realize Gibbs and the KMS condition. Null Landauer corresponds
to bijective transport of the complete fiber; preparation, projection, and
reset receive separate balances when the coupling discards a
fiber. The thermal reading first constructs the internal transducer and then applies
the unit chart.

Bilayer responses produce directly impedance, velocity,
permeability and permittivity through
\[
 \widehat Z=\sqrt{\widehat\mu/\widehat\varepsilon},
 \qquad
 \widehat c=(\widehat\mu\widehat\varepsilon)^{-1/2}.
\]
The same causal precedence governs the electromagnetic reading: the internal
constitutive operator precedes the realization of
\(\varepsilon_0,\mu_0,Z_0\) and \(c\) in the International System of Units.

The gravitational branch first constructs, from the outputs of the enriched
state, the action scale, the gravitational selector and the section
\(G_a\). The Palatini--Cartan--Holst variational formulation then receives
these outputs, together with \(\gamma_{\rm HMT}\), and admits the constrained
Hamiltonian transform
\[
 H_{\rm tot}=\int_\Sigma
 (\lambda^i G_i+N^a V_a+N\mathcal H).
\]
The subsequent equivalence between covariant action, constrained Hamiltonian system and
Einstein--Schrödinger evolution preserves co-tetrad, connection, curvature and
torsion within the same scheme. The torsion-free collective sector carries
the gravitational excitation; local torsion induced by spin density
supplies the material channel. The type, transducer and chart of \(G\) precede
this dynamical realization and are not selected through it.

The solenoidal and gauge branches receive realizations of the common finite operator
\[
 \mathcal L_{\rm cyc}^{\rm TPK}=2I-C-C^{-1}=3(I-P_0)=3I-J,
\]
\[
 \mathfrak R_{\rm sol}(\mathcal L_{\rm cyc}^{\rm TPK})
 =\mathcal L_{\rm H}^{(0)},\qquad
 \mathfrak R_{\rm gau}(\mathcal L_{\rm cyc}^{\rm TPK})
 =\mathcal L_{\rm H}^{(1)}.
\]
The hydrodynamic branch transports it through the Galerkin--register intertwiner, with coercivity and residual declared; the gauge branch transports it through the gauge functional and the transfer--gap operator. The finite proof and the bridges to the continuous equations are presented with their distinct hypotheses. The self-inertial frames, the Machian factorization, and sheet cosmology are developed after this local and collective construction, where they acquire a domain, a reader, and a typed physical interpretation.

\subsection*{Method of proof and contrast}

Each result is presented through its object, domain, map, derivation,
premises, certificate, and scope. Reproducible calculations accompany the
proof and do not replace it. Metrological charts are applied after
the internal scalar or operator has been constructed. CODATA and PDG comparisons
are opened after the HMT output has been sealed\@. This order preserves the
causal direction of the program: the structure generates the digit and the magnitude;
the external reference measures the agreement attained by HMT\@.

The work is organized into five causally ordered Parts. The first four
construct the formal core, the discrete continuum, the genealogy of the
constants, and the five specialization chains; the fifth realizes that architecture
in magnitudes, particles, fields, and physical geometry. This sequence expresses
the mathematical dependence among structure, evaluation, closure, and dimensional
realization.

\subsection*{Composition of the treatise and scope of the statements}

Part I constructs APP from the nonadic cell, presents TRIT as an
oriented lift, and defines the Topological Phase Kernel with its fibers,
transports, memory, clock, joint nonadic connection, and compatible
prolongation. The horizon line opens this construction through the
notions of scale, boundary, and measure.

Part II develops the discrete structure of the continuum. The double
projection separates Archimedean evaluation from incidential publication.
From Paley--Witt, the Steiner--Golay design branch and the
Niemeier--Leech--\(V^\natural\)--Monster--Moonshine lattice branch are distinguished. Within that same structure
appear the internal square root of \(-1\), the three quadratic realizations, and the
spectral--polar, solenoidal, gauge, cohomological, and
determinantal substructures, typed by their own maps and terminals.

Part III establishes the genealogy of the fundamental constants: the
correlated coinductive publication of
\(\pi,e,\varphi,\alpha_{\rm HMT}\), the specific readers and the two internal
pathways of \(\alpha_{\rm HMT}\), \(-1/12\), the action scale, Planck,
Boltzmann, the Barbero--Immirzi parameter, CKM--CP, the vacuum
constitutives, the gravitational constant and the atlas of mathematical and
physical constants.

\Needspace{12\baselineskip}
Part IV applies the five substructures constructed in Part II in
chapters of identical editorial rank. Each chapter sets out the internal
result, the specialized comparator, its premises, and the exact criterion of
publication. The conclusion in the corresponding conventional formulation is
deduced only when the comparator identity has been verified throughout its
domain.

Part V realizes APP--TRIT--TPK through the physical readers: torus and sheets,
temporal semantics, curvatures, channels, and holonomy. On that grammar it
defines the particle-route, the bidirectional electron, the mass families,
the constitutive laws, statistics, information, mixing, variational gravity,
self-inertia, sheet cosmology, and the interface with M-theory.

The appendices gather auxiliary proofs, complete censuses, seeds, decimal strings, extensive tables, auxiliary derivations, and reproducible calculations. The main body preserves the sequence of definition, construction, theorem, interpretation, and result. Each claim distinguishes its mathematical domain, its starting structure, its realization chart, and its external test. This separation permits the exact closures, physical realizations, and extension programs to be presented together while preserving for each the degree of proof proper to it.
